\scanpage{14}
\diagram{\begin{tikzcd}
H^2(\mathbf Q,H) \arrow[d,equal] \\
\varinjlim H^2(G,K^*\otimes\check M)
\end{tikzcd}}

\diagram{\begin{tikzcd}[column sep=small,row sep=large]
H^1(G,\unsure{J|D}\otimes\check M) \arrow[r,dash] \arrow[d,equal] & H^1(G,C\otimes\check M) \arrow[r] \arrow[d,equal] & H^2(G,K^*\otimes\check M) \arrow[r] \arrow[d,equal] & H^2(G,J\otimes\check M) \arrow[r,dash] \arrow[d,equal] & H^2(G,C\otimes\check M) \arrow[d,equal] \\
\sum_p\widehat H^{-1}(G,D_p\otimes\check M) \arrow[r] \arrow[d,equal] & \widehat H^{-1}(G,\check M) \arrow[r] \arrow[d,equal] & \widehat H^0(G,\check M\otimes\unsure{Y|X}) \arrow[r] & \sum_p\underbrace{\widehat H^0(G,D_p\otimes\check M)}_{\widehat H^0(G_p,\check M)} \arrow[r] \arrow[d,equal] & \widehat H^0(G,\check M) \\
\displaystyle\sum_p\underbrace{\widehat H^{-1}(G_p,\check M)}_{\begin{array}{c}=\\\underbrace{\operatorname{Ker}(\check M_{G_p}\xrightarrow{N_{G_p}}\check M^{G_p})}_{\begin{array}{c}\text{pour }p=\infty,\text{ c'est }\check M/2\check M\\\text{dans le cas compact}\end{array}}\ \deleted{\unsure{\check M}}\end{array}}
& \underbrace{\operatorname{Ker}(\check M_G\xrightarrow{N_G}\check M^G)}_{\begin{array}{c}\unsure{\check M_G}\\\text{quotient de }\check M/2\check M\end{array}} & & \check M^{G_p}/N_{G_p}\check M \arrow[d,equal] & \\
& \text{\unsure{identiq\textsuperscript{t}} nul si }H\text{ compact} \arrow[u,bend left=20,no head] & & H^2(\mathbf Q_p,H^0)\simeq H^2(\mathbf Q_p,H/\mu_2) &
\end{tikzcd}}

\diagram{\begin{tikzcd}
\unsure{0} \arrow[r] & H/\mu_2 \arrow[r] & G/\mu_2 \arrow[r] & G' \arrow[r] \arrow[d,equal] & 0 \\
& & & G_m\times\widehat{\mathbf Z} &
\end{tikzcd}}

\diagram{\begin{tikzcd}
\Gamma^0/(\text{racines de }1) \arrow[r,phantom,"\simeq"] & M \arrow[r] & \unsure{\mathbf Z|\mathbf Z^+} & \text{\uline{stable par }}G_p\ \deleted{\text{décomposition }\unsure{\text{(du groupe)}}\ \text{\illegible{2}}} \\
& & \begin{array}{l}[\text{valuation }\deleted{\text{p-adique}}\\\text{p-adique normalisée par }v_p(p)=1]\end{array} \arrow[u,no head,bend left=25] &
\end{tikzcd}}

$H^0$ est dual du groupe $M$ formé des \deleted{\unsure{nb}} nbs de Weil qui sont
\[
\deleted{\unsure{\mathfrak p}}\qquad
\unsure{\mathfrak p_1|\rho_1}\quad\overline{\mathfrak p}_1\quad\mathfrak p_2\quad\overline{\mathfrak p}_2\quad\cdots\quad\mathfrak p_\nu,\overline{\mathfrak p}_\nu
\]
\[\unsure{n_1}\]

\diagram{\begin{tikzcd}
& \Gamma^0=\operatorname{Ker}i^*:\ \text{\uline{\inserted{p-}unités de Weil}} \arrow[d] & & \\
\mathbf Z \arrow[r,"\epsilon^*"] & \Gamma \arrow[r,"i^*"] & \mathbf Z & i^*\epsilon^*=2\operatorname{Id} \\
& \text{nbs de Weil} & &
\end{tikzcd}}
\[
\left\{\begin{array}{l}
\epsilon^*(1)=p\\
|\lambda|=p^{i^*(\lambda)/2}
\end{array}\right.
\]
\diagram{\begin{tikzcd}
\deleted{\mathbf Z} \arrow[r,phantom,"{\deleted{\longrightarrow}}"{description}] & \deleted{\mathbf Z/2\mathbf Z} \arrow[r,phantom,"{\deleted{\longrightarrow}}"{description}] & \deleted{\Gamma/\epsilon^*(\mathbf Z)}
\end{tikzcd}}

\diagram{\begin{tikzcd}[row sep=large]
& 0 & \\
& \mathbf Z/2\mathbf Z\ \deleted{\unsure{p-p^{1/2}}} \arrow[u,dash] & \\
\begin{array}{l}\text{en envoyant}\\1\text{ sur la}\\\text{classe de }\unsure{\pm p^{1/2}}\end{array} \arrow[r,bend right=20] & \Gamma/(\text{racines de }1).\epsilon^*(\mathbf Z) \arrow[u] \arrow[r,leftrightarrow] & U \\
& \Gamma^0/\text{racines de }1 \arrow[u] \arrow[r,leftrightarrow] & U^0 \\
& 0 \arrow[u,dash] &
\end{tikzcd}}

\margin{$u=\epsilon p^\nu$\quad $\epsilon$ racine de $1$\quad $\nu\in\mathbf Z$\par
$u$ unité de Weil\par
$\Rightarrow\nu=0$, donc $u$ racine de $1$}
\mspage{13}
